% One-time imported draft. Editable native source; no Markdown regeneration.
\section{Normalize a value with another instruction
  sequence}\label{004-bit-inspector-p10-normalize-a-value-with-another-instruction-sequence}

The alternative loads the signed word, shifts its low word left by 32
bits, then shifts logically right by 32 bits. The first operation
discards the former high half; the second fills the new high half with
zeroes. The final value is the same zero-extended word that \texttt{lwu}
produces here.

It also computes the sum at doubleword width and then applies a word
shift by zero. Although no low-word bit position moves, that word
instruction still sign-extends its 32-bit result. It therefore
normalizes the sum to the required word-width result.

A zero shift is not necessarily a full-register identity operation. The
familiar \texttt{nop} encoding uses the zero register as both source and
destination, so its word semantics produce no useful state change.
Applying a word operation to a nonzero register can change the upper
half even when its shift distance is zero.
